\chapter*{Declaración de uso de inteligencia artificial}
\addcontentsline{toc}{chapter}{Declaración de uso de inteligencia artificial}

Esta declaración recoge el uso que he hecho de herramientas de inteligencia artificial generativa durante el trabajo. La incluyo al principio, por transparencia con quien lo lea: qué herramientas y modelos usé y en qué periodos, para qué tareas, y qué parte del trabajo es mía y cómo la verifiqué.

\section*{Herramientas y modelos}

\begin{itemize}
    \item \textbf{Claude Code (Anthropic).} Ha sido la herramienta principal~\citep{claudecode} desde enero de 2026 hasta la entrega, en septiembre de 2026. Es el asistente de programación de Anthropic para el terminal y el editor, y trabaja directamente sobre el repositorio del proyecto. En cada momento usé la versión de la familia Claude Opus disponible en la herramienta y, desde junio de 2026, también la familia Claude Fable. Con los identificadores de la documentación de Anthropic~\citep{anthropic2026models}, los modelos publicados en ese periodo fueron Claude Opus 4.5 (\texttt{claude-opus-4-5-20251101}, noviembre de 2025, el vigente en enero de 2026), Claude Opus 4.6 (\texttt{claude-opus-4-6}, febrero de 2026), Claude Opus 4.7 (\texttt{claude-opus-4-7}, abril de 2026), Claude Opus 4.8 (\texttt{claude-opus-4-8}, mayo de 2026), Claude Fable 5 (\texttt{claude-fable-5}, junio de 2026), Claude Opus 5 (\texttt{claude-opus-5}, julio de 2026) y Claude Fable 5.1 (\texttt{claude-fable-5-1}, septiembre de 2026). No conservo un registro de qué modelo atendió cada sesión, así que declaro el rango completo.
    \item \textbf{Figma desde Claude Code.} El prototipo navegable de Figma se construyó y se actualizó desde Claude Code a través del servidor MCP oficial de Figma~\citep{figmamcp} (\textit{Model Context Protocol}), que permite al asistente leer y modificar el fichero de diseño. Las pantallas del segundo prototipo y su actualización tras la auditoría entre diseño e implementación (capítulos~\ref{cap:analisisRequisitos} y~\ref{cap:disenoSistema}) se hicieron así, sobre mis indicaciones y con mi revisión en Figma.
    \item \textbf{Gemini (Google).} Antes de Claude Code, en junio y julio de 2025 y en diciembre de 2025, usé la aplicación Gemini de Google con los modelos que ofrecía entonces: en junio y julio de 2025 la familia Gemini 2.5 (Gemini 2.5 Pro y Gemini 2.5 Flash, publicados como versiones estables ese mismo junio) y en diciembre de 2025 la familia Gemini 3 (Gemini 3 Pro y Gemini 3 Flash, publicados en noviembre y diciembre de 2025) junto con la 2.5~\citep{google2026gemini}. No conservo registro de qué variante respondió a cada consulta.
\end{itemize}

No he usado ChatGPT ni ninguna otra herramienta de inteligencia artificial generativa.

\section*{Tareas en las que se usaron}

Gemini lo usé para desarrollar la idea y como apoyo en los primeros pasos del proyecto, sin una tarea concreta.

Claude Code fue mi asistente de forma recurrente en el resto del proyecto. Implementé el código del repositorio con su ayuda: le daba mis especificaciones, revisaba y corregía lo que proponía y lo verificaba con mis bancos de pruebas. Diseñé las pantallas de Figma con el asistente trabajando sobre mis indicaciones, preparé con él los experimentos del capítulo~\ref{cap:evaluacion} y los scripts que producen sus tablas y figuras, y lo usé también en la documentación interna del proyecto. Esta memoria la escribí yo, y el asistente me sirvió para perfeccionarla y mejorarla: le pedía borradores a partir de mis indicaciones y de la documentación del proyecto, y los revisé, corregí y reescribí hasta dejarlos en mi estilo; también me ayudó con la traducción al inglés.

\section*{Autoría y verificación}

El sistema lo implementé yo y la memoria la escribí yo, con la inteligencia artificial como asistente en los dos casos. Las decisiones de alcance, de arquitectura y de modelado son mías, tomadas en las reuniones con los directores del TFG y a partir de lo que pedían los nutricionistas consultados. El código que el asistente proponía lo revisé y lo corregí, y ninguna pieza entró en el sistema desplegado sin pasar por los bancos de pruebas que definí para cada bloque. Las entrevistas con profesionales y la prueba con usuarios las hice yo. Los datos y cifras de la evaluación los medí sobre el sistema; no los generó ningún modelo. La bibliografía la verifiqué en su fuente. Y ninguna frase de esta memoria se ha incorporado sin mi revisión.
